\documentclass[12pt,a4paper]{article}

% ------------------------------------------------
% Pacotes
% ------------------------------------------------
\usepackage[utf8]{inputenc}
\usepackage[T1]{fontenc}
\usepackage[brazil]{babel}
\usepackage{graphicx}
\usepackage[table]{xcolor}
\usepackage{array}
\usepackage{tabularx}
\usepackage{geometry}
\usepackage{fancyhdr} 
\usepackage{mathptmx} % Times New Roman
\usepackage{enumitem} 
\usepackage{amsfonts}
\usepackage{amsmath}
\usepackage{amssymb}
\usepackage{tikz}
\usetikzlibrary{positioning, shapes.geometric}

% TikZ Styles for ER Diagrams
\tikzset{
    block/.style={rectangle, draw, thick, minimum width=2cm, minimum height=1cm, fill=white},
    diamond_node/.style={diamond, draw, thick, minimum width=2cm, minimum height=1cm, fill=white, aspect=2},
    attribute/.style={ellipse, draw, thin, minimum width=1.5cm, minimum height=0.6cm, fill=white, font=\scriptsize},
    line/.style={draw, thick},
    specialization/.style={draw, circle, inner sep=1pt, font=\tiny, fill=white}
}

% ------------------------------------------------ 
% Configuração de página
% ------------------------------------------------
\geometry{
    a4paper,
    left=2cm,
    right=2cm,
    top=2cm,
    bottom=2cm,
    headheight=15pt
}

\pagestyle{fancy}
\fancyhf{}
\renewcommand{\headrulewidth}{0pt}
\setlength{\parindent}{0pt}

% Caixa para resposta
\newcommand{\resposta}[1]{
\noindent
\framebox[\textwidth]{
\begin{minipage}{0.98\textwidth}
\vspace{#1}
\end{minipage}
}
}

\newcolumntype{L}[1]{>{\raggedright\arraybackslash}p{#1}}

% ------------------------------------------------
% Documento
% ------------------------------------------------
\begin{document}

% ------------------------------------------------
% Cabeçalho
% ------------------------------------------------
\begin{tabularx}{\textwidth}{@{} L{4cm} X @{}}
    \raisebox{-0.5cm}{\includegraphics[width=4cm]{Universidade_de_Rio_Verde.jpg}} &
    \hfill
    \begin{tabular}{|>{\centering\arraybackslash}m{10cm}|>{\raggedright\arraybackslash}p{1.8cm}|}
        \hline
        \cellcolor{black}\color{white}\rule{0pt}{1.2cm}\textbf{ATIVIDADE PRÁTICA: MODELAGEM ER E LÓGICA} 
        & 
        \scriptsize Nota \tabularnewline
        \hline
    \end{tabular} 
\end{tabularx}
\vspace{0.3cm}

\begin{tabularx}{\textwidth}{|X|L{2cm}|}
\hline
\textbf{Curso} & \textbf{Campus} \\
\textit{Engenharia de Software} & \textit{Rio Verde} \\
\hline
\multicolumn{2}{|l|}{\textbf{Disciplina}} \\  
\multicolumn{2}{|l|}{ \textit{Banco de Dados I}} \\
\hline
\textbf{Nome do(a) acadêmico(a)} & \textbf{Valor} \\
& 10,0 \\
\hline
\textbf{Professor} & \textbf{Turma} \\
\textit{ Sergio Souza Novak} & \textit{A/B} \\
\hline
\textbf{Data} & \textbf{Código} \\
\textit{17/03/2026} & \textit{ESW} \\
\hline
\end{tabularx}

\vspace{0.5cm}

\begin{center}
    \rule{\textwidth}{0.4pt}
    \textbf{ENUNCIADO DO CENÁRIO: SISTEMA CLINICARV}
\end{center}

A Clínica Médica "ClinicaRV" necessita de um novo sistema de banco de dados para gerenciar seus atendimentos. Após o levantamento de requisitos, foram identificadas as seguintes regras de negócio:

\begin{enumerate}
    \item \textbf{Médicos e Especialidades:} Todo médico possui um CRM (único), nome e telefone de contato. Cada médico deve estar vinculado a exatamente uma especialidade (ex: Cardiologia, Pediatria). Uma especialidade pode ter vários médicos cadastrados.
    \item \textbf{Pacientes e Atendimento:} Os pacientes são identificados pelo CPF e possuem nome e data de nascimento. Um paciente pode realizar várias consultas na clínica ao longo do tempo.
    \item \textbf{Convênios:} A clínica atende por diversos convênios (ex: Unimed, Ipasgo). Um registro de convênio contém um código identificador, o nome do plano e a empresa responsável. Um paciente pode utilizar apenas um convênio por vez (ou nenhum), enquanto um convênio atende a milhares de pacientes. 
    \item \textbf{Credenciamento:} Os médicos podem atender por vários convênios diferentes, e um convênio pode ter diversos médicos credenciados para atendimento.
    \item \textbf{Consultas:} Uma consulta registra a data, o horário e o valor cobrado. Toda consulta deve estar obrigatoriamente vinculada a um médico e a um paciente.
    \item \textbf{Exames:} Durante uma consulta, o médico pode solicitar vários exames (ex: Hemograma, Raio-X). Cada exame solicitado possui um código, nome e preço. Um exame solicitado está vinculado diretamente à consulta que o originou.
    \item \textbf{Receitas e Medicamentos:} Cada consulta gera obrigatoriamente uma única receita médica (prescrição), identificada por um número sequencial e contendo observações gerais do médico.
    \item \textbf{Prescrição de Medicamentos:} Uma receita pode conter vários medicamentos diferentes, e o mesmo medicamento (identificado por código e nome comercial) pode ser prescrito em diversas receitas de pacientes distintos.
\end{enumerate}

\vspace{0.8cm}

\begin{center}
    \rule{\textwidth}{0.4pt}
    \textbf{QUESTÕES}
\end{center}

\begin{enumerate}[label=\textbf{Questão \arabic*:}, leftmargin=*]
    \item \textbf{Modelo Entidade-Relacionamento (Conceitual):} 
    Com base no cenário acima, desenhe o Diagrama Entidade-Relacionamento (DER). Certifique-se de incluir:
    \begin{itemize}
        \item Todas as 8 entidades citadas;
        \item Os atributos principais (incluindo chaves primárias sublinhadas);
        \item Os relacionamentos e suas respectivas cardinalidades (mínima e máxima).
    \end{itemize}
    
    \resposta{12cm}

    \newpage

    \item \textbf{Mapeamento para o Modelo Lógico (Relacional):}
    Converta o DER elaborado na questão anterior para o Modelo Lógico Relacional. Apresente o esquema das tabelas seguindo o padrão: \textbf{TABELA}(\underline{PK}, Atributo, FK\_Referencia). 
    
    \textit{Dica: Atenção redobrada aos relacionamentos N:N e para onde as Chaves Estrangeiras (FK) devem ser propagadas nos relacionamentos 1:N.}

    \resposta{15cm}

\end{enumerate}

\end{document}
